\documentclass[11pt,a4paper]{article}
\usepackage[margin=1in]{geometry}
\usepackage[T1]{fontenc}
\usepackage{lmodern}
\usepackage{amsmath,amssymb,amsthm,bm}
\usepackage{booktabs}
\usepackage{microtype}
\usepackage[hidelinks]{hyperref}

\newtheorem{assumption}{Assumption}
\newtheorem{proposition}{Proposition}
\newtheorem{remark}{Remark}
\newcommand{\E}{\mathbb{E}}
\newcommand{\Var}{\operatorname{Var}}
\newcommand{\Cov}{\operatorname{Cov}}
\newcommand{\tr}{\operatorname{tr}}
\newcommand{\CN}{\mathcal{CN}}

\title{Choosing a Spectral Prediction Target for EEG:\\
An Idealized Analysis of Five Classical Frequency Bands}
\author{}
\date{}

\begin{document}
\maketitle

\begin{abstract}
This note develops a theoretical rationale for masked prediction of EEG band power rather
than sample-wise waveform reconstruction. The target is defined using five \emph{ideal,
non-overlapping classical frequency bands}: $\delta$ (0.5--4 Hz), $\theta$ (4--8 Hz),
$\alpha$ (8--13 Hz), $\beta$ (13--30 Hz), and $\gamma$ (30--45 Hz). In an idealized
Gaussian spectral model, Fourier phases are ancillary for a state that changes only the
power spectrum. If the spectrum is constant within each band, the five band-power
statistics are sufficient for the five spectral levels; averaging within a band reduces
periodogram noise, and the logarithm makes the target's variance independent of those
levels. A conditional-risk decomposition then identifies the state-dependent and
irreducible components of masked prediction. These are statements about an explicitly
specified model, not a proof that a particular EEG encoder or training objective must
outperform waveform reconstruction in practice.
\end{abstract}

\section{Question and target}

A masked EEG encoder receives a context $C$ containing visible channel--time patches and
predicts a property of a masked patch $x_M$. Which property should it predict? For a
downstream label that depends primarily on sustained rhythmic activity, a natural
candidate is the patch's spectral energy in a small number of physiologically meaningful
bands. This choice removes the requirement to reproduce the precise sample-by-sample
realization of the masked waveform. It does \emph{not} assert that all EEG labels are
determined by band power: transient morphology, spatial relationships, and phase coupling
can also be essential.

\subsection{Why power? Rhythmic strength without exact phase}

Consider first a single oscillation,
$x(t)=A\cos(2\pi ft+\phi)$. Its mean square over a window containing an integer
number $q\geq1$ of periods, $T=q/f$, is
\begin{align}
\frac{1}{T}\int_0^T x(t)^2\,dt
&=\frac{A^2}{T}\int_0^T\cos^2(2\pi ft+\phi)\,dt\nonumber\\
&=\frac{A^2}{2}
 +\frac{A^2}{2T}\int_0^T\cos(4\pi ft+2\phi)\,dt
 =\frac{A^2}{2}.
\label{eq:sinusoidpower}
\end{align}
Thus mean-square power measures the oscillation's squared amplitude but does not
depend on where its peak happens to fall within a complete-cycle window.
Squaring alone does \emph{not} remove phase: $x(t)^2$ still oscillates and still
depends on $\phi$. Phase invariance comes from averaging over complete periods,
or, for an ideal frequency-domain target, from taking squared Fourier magnitudes.
For a finite window containing a non-integer number of periods, the boundary
term need not vanish. This matters especially for slow rhythms in short patches.

For a multi-frequency signal, ideal band-pass filtering isolates a classical
rhythm range. Parseval's identity equates the filtered signal's mean square
with the spectral energy in that range; the precise finite-DFT normalization
appears in Eq.~\eqref{eq:parseval}. This is why the power in, for example,
the $\alpha$ or $\delta$ band is a natural measure of that rhythm's strength.
Power does not preserve the phase relations among frequency components or
channels, which may themselves carry information for other tasks.

The quantity expected to persist across neighbouring patches is the
\emph{state-dependent spectral level}, not each patch's exact measured power.
If a state lasts longer than one patch, visible patches can provide information
about the expected energy of a masked patch. Its realized power still fluctuates,
and a stable oscillator may have partly predictable phase across adjacent
patches. The formal shared-state and conditional-independence conditions needed
for the masked-prediction argument are stated separately in
Assumption~\ref{as:context}; they are not consequences of Eq.~\eqref{eq:sinusoidpower}.

\subsection{Why take the logarithm?}

Power is nonnegative and often spans a wide range. For positive values, a
logarithm compresses this range and makes a difference in log power approximately
a \emph{relative} difference in power:
\begin{equation}
\log(P+\Delta P)-\log P\approx\frac{\Delta P}{P}
\qquad (|\Delta P|\ll P).
\label{eq:relativechange}
\end{equation}
If amplitude is multiplied by $g>0$, power is multiplied by $g^2$, and log
power is shifted by $2\log g$. This turns a multiplicative factor into an
additive one; it does not by itself remove the factor. A fixed reference power
or normalization makes the argument of the logarithm dimensionless.

There is also a statistical reason. Individual periodogram powers fluctuate
roughly in proportion to their spectral level. Under the ideal model developed
below, band-averaged power has conditional variance proportional to the
\emph{square} of that level, whereas log band power has conditional variance
that depends only on the number of frequency bins (Proposition~\ref{prop:log}).
Thus logarithmic power is a natural target when squared error should not be
dominated solely by the bands or states with the largest absolute power.
This argument applies exactly to $\log P$ in that model. An offset transform
such as $\log(a+P)$ behaves differently at low power, as
Eq.~\eqref{eq:logoffset} will show.

We adopt the following conventional, half-open positive-frequency bands throughout:
\begin{equation}
\begin{aligned}
B_{\delta}&=[0.5,4), & B_{\theta}&=[4,8), & B_{\alpha}&=[8,13),\\
B_{\beta}&=[13,30), & B_{\gamma}&=[30,45]\quad\text{Hz}.
\end{aligned}
\label{eq:bands}
\end{equation}
The exact boundary convention is stipulated for the analysis; physiological rhythms do
not have sharp boundaries. In particular, the word ``classical'' refers to these five
\emph{ideal rectangular bands}, not to five generic or overlapping spectral filters.

For one real-valued, mean-zero patch $x=(x_0,\ldots,x_{N-1})$ sampled at $f_s$, let
\begin{equation}
X_k=\sum_{n=0}^{N-1}x_n e^{-2\pi i kn/N},\qquad
I_k=\frac{|X_k|^2}{N},\qquad f_k=\frac{kf_s}{N}.
\label{eq:dft}
\end{equation}
We use positive-frequency bins $1\leq k\leq\lfloor(N-1)/2\rfloor$, excluding DC and,
when present, Nyquist. For $\mathcal K_b=\{k:f_k\in B_b\}$ and
$m_b=|\mathcal K_b|>0$, define the \emph{mean band periodogram}
\begin{equation}
P_b=\frac{1}{m_b}\sum_{k\in\mathcal K_b}I_k,
\qquad Y_b=\log P_b.
\label{eq:bandpower}
\end{equation}
The vector $(Y_{\delta},Y_{\theta},Y_{\alpha},Y_{\beta},Y_{\gamma})$ is the
idealized five-dimensional prediction target. Since the Gaussian model below has
$P_b>0$ almost surely, its logarithm is well-defined once power is expressed
relative to a fixed reference unit. Such a reference adds only a fixed offset
to $Y_b$ and does not alter the results below.

For illustration, at $N=200$ and $f_s=256$ Hz, the spacing is $\Delta f=1.28$ Hz.
The number of represented positive-frequency bins is small in the slow bands:
\begin{table}[h]
\centering
\caption{Classical bands and represented DFT bins for a 200-sample patch at 256 Hz.
The last column is the theoretical standard deviation of $\log P_b$ under
Assumption~\ref{as:spectral}; it is not an empirical EEG measurement.}
\label{tab:bins}
\begin{tabular}{lccc}
\toprule
Band & Range (Hz) & $m_b$ & $\sqrt{\psi_1(m_b)}$ \\
\midrule
$\delta$ & $[0.5,4)$ & 3 & 0.63 \\
$\theta$ & $[4,8)$ & 3 & 0.63 \\
$\alpha$ & $[8,13)$ & 4 & 0.53 \\
$\beta$ & $[13,30)$ & 13 & 0.28 \\
$\gamma$ & $[30,45]$ & 12 & 0.29 \\
\bottomrule
\end{tabular}
\end{table}
There is no bin between 0.5 and 1.28 Hz in this example. Calling the first statistic
``$\delta$ power'' therefore does not imply that a short patch resolves the entire
0.5--4 Hz band. Longer windows or a different estimator are required for finer
low-frequency resolution.

\section{An idealized spectral model}

The following finite-dimensional model is chosen so that its conclusions are exact
within the stated observation space. For a weakly stationary Gaussian process on a
long window, a Fourier-diagonal likelihood is instead an approximation, commonly
associated with the Whittle likelihood. Nothing below assumes that real EEG is
exactly stationary or Gaussian.

\begin{assumption}[Five-band Gaussian spectral state]\label{as:spectral}
Let $\Theta$ denote a latent state. Conditional on $\Theta=\theta$, the positive-frequency
Fourier coefficients of a real, mean-zero patch in the five bands are independent and
\begin{equation}
X_k\mid\theta\sim\CN\!\left(0,Ns_b(\theta)\right)
\quad\text{for }k\in\mathcal K_b,\qquad s_b(\theta)>0.
\label{eq:complexnormal}
\end{equation}
The negative-frequency coefficients are their complex conjugates. The observation for
the sufficiency statement consists only of coefficients in
$\bigcup_b\mathcal K_b$; any omitted frequencies are either unobserved or have a
distribution independent of $\theta$. The state affects this observation only through
$\bm s(\theta)=(s_{\delta},s_{\theta},s_{\alpha},s_{\beta},s_{\gamma})$.
\end{assumption}

This is a \emph{piecewise-flat, ideal rectangular-band} model. In a continuous-time
interpretation, $s_b$ describes the spectral level within band $B_b$. An actual
spectrum may vary within a band and may change smoothly at its boundaries. In that
case, the exact sufficiency result below no longer follows.

\subsection{Phase carries no state information in the model}

Write $X_k=R_ke^{i\Phi_k}$. Circular symmetry of the complex normal in
Eq.~\eqref{eq:complexnormal} gives
\begin{equation}
I_k\mid\theta\sim s_b(\theta)\,\mathrm{Exp}(1),\qquad
\Phi_k\mid\theta\sim\mathrm{Unif}[0,2\pi),
\label{eq:exponential}
\end{equation}
with independent magnitudes and phases. The phase distribution does not depend on
$\theta$.

\begin{proposition}[Ancillarity of phase]\label{prop:phase}
Under Assumption~\ref{as:spectral}, the phase vector of the observed positive-frequency
coefficients is ancillary for $\bm s(\theta)$. The likelihood for $\theta$ depends on
the observation only through the periodogram values $\{I_k\}$.
\end{proposition}
\begin{proof}
Each conditional complex-normal density is proportional to
$s_b(\theta)^{-1}\exp[-I_k/s_b(\theta)]$. In polar coordinates, its angular factor
is constant. Multiplication across the observed bins leaves no phase term that
depends on $\theta$.
\end{proof}

This statement does \emph{not} say that Fourier phase is always useless in EEG.
It says that phase is uninformative about a state \emph{defined to act only through
the five spectral levels} in this model. Nor does it say that every aspect of
amplitude is predictable: each $I_k$ retains an exponential fluctuation even if
$\theta$ is known.

\subsection{Why band averaging is sufficient under piecewise flatness}

\begin{proposition}[Band-power sufficiency and estimation]\label{prop:sufficient}
Under Assumption~\ref{as:spectral}, $(P_b)_{b\in\{\delta,\theta,\alpha,\beta,\gamma\}}$
is sufficient for $\bm s$ in the observed, band-limited Fourier data. For each band,
the maximum-likelihood estimate is $\widehat{s}_b=P_b$; by invariance of maximum
likelihood, $\log P_b$ estimates $\log s_b$.
\end{proposition}
\begin{proof}
Ignoring terms independent of $\bm s$, the negative log-likelihood is
\begin{equation}
-\log p(\{X_k\}\mid\bm s)
=\sum_b\left(m_b\log s_b+\frac{m_bP_b}{s_b}\right)+\mathrm{const}.
\label{eq:likelihood}
\end{equation}
The factorization criterion gives sufficiency. Differentiating the $b$th term
with respect to $s_b$ yields
$m_b/s_b-m_bP_b/s_b^2=0$, so its positive minimizer is $s_b=P_b$.
\end{proof}

The frequency-domain mean $P_b$ is proportional to the time-domain mean square
of an \emph{ideal circular band-pass} signal. With the DFT normalization in
Eq.~\eqref{eq:dft}, Parseval's identity gives
\begin{equation}
\frac{1}{N}\sum_{n=0}^{N-1}|x^{(b)}_n|^2
=\frac{2}{N^2}\sum_{k\in\mathcal K_b}|X_k|^2
=\frac{2m_b}{N}P_b.
\label{eq:parseval}
\end{equation}
The two band-power conventions differ only by the known factor $2m_b/N$.
Their logarithms differ by a band-dependent constant. This equivalence presumes
the stated ideal frequency mask and excludes DC and Nyquist.

\begin{remark}
Sufficiency holds for the five parameters $\bm s$ \emph{within the model and the
specified observation space}. If the spectral shape inside a band, the power outside
the five bands, cross-channel relationships, or transient morphology also varies
with the state, five band averages need not retain all state information.
\end{remark}

\subsection{Log power and estimator noise}

\begin{proposition}[Distribution and variance stabilization]\label{prop:log}
Under Assumption~\ref{as:spectral},
\begin{equation}
P_b=s_bG_b,\qquad G_b\sim\mathrm{Gamma}(m_b,\text{scale}=1/m_b).
\label{eq:gamma}
\end{equation}
Writing $\psi$ and $\psi_1$ for the digamma and trigamma functions, respectively,
\begin{align}
\E[\log P_b\mid\theta]
&=\log s_b(\theta)+\psi(m_b)-\log m_b,\label{eq:logmean}\\
\Var(\log P_b\mid\theta)
&=\psi_1(m_b).\label{eq:logvar}
\end{align}
Thus the conditional noise variance on the log scale depends on the number of
frequency bins, but not on the state's spectral level.
\end{proposition}
\begin{proof}
The sum of $m_b$ independent unit-mean exponentials is Gamma with shape $m_b$
and unit scale. Dividing by $m_b$ gives Eq.~\eqref{eq:gamma}. The standard log-Gamma
moments are $\E\log G_b=\psi(m_b)-\log m_b$ and
$\Var(\log G_b)=\psi_1(m_b)$.
\end{proof}

For comparison, a single periodogram ordinate has
$\operatorname{sd}(\log I_k\mid\theta)=\pi/\sqrt{6}\approx1.28$.
Band averaging lowers this fluctuation as $m_b$ grows, with
$\psi_1(m_b)\approx1/m_b$ for large $m_b$. The benefit is limited in the
slow bands of short patches, where $m_b$ is only three in Table~\ref{tab:bins}.
The transformed target is not noise-free: $G_b$ is a genuine stochastic
realization that remains after conditioning on the state.

An additive offset inside the logarithm changes the statement. For the
mathematical family $Y_b^{(a)}=\log(a+P_b)$ with $a>0$, a first-order delta-method
approximation around $G_b=1$ gives
\begin{equation}
\Var(Y_b^{(a)}\mid\theta)
\approx\frac{1}{m_b}\left(\frac{s_b(\theta)}{a+s_b(\theta)}\right)^2.
\label{eq:logoffset}
\end{equation}
Unlike Eq.~\eqref{eq:logvar}, this variance depends on the state unless the
spectral level is large relative to $a$. Equation~\eqref{eq:logoffset} is an
approximation about a mathematical alternative, not part of the exact
sufficiency argument.

\section{Masked prediction and irreducible uncertainty}

The sufficient-statistic result concerns inference from an \emph{observed} patch.
Masked prediction asks a different question: what can a model infer about a
\emph{missing} patch from its context? A separate temporal assumption is required.

\begin{assumption}[Shared state and conditional independence]\label{as:context}
The visible context $C$ and the masked target $Y$ are associated with a latent
state $\Theta$. Conditional on $\Theta$, $Y$ is independent of $C$. The target
has a finite second moment.
\end{assumption}

This is an idealization of state persistence across patches. It is not implied
by Assumption~\ref{as:spectral}; adjacent samples of a continuous process may
remain correlated even after its broad state is fixed.

\begin{proposition}[Bayes risk decomposition]\label{prop:risk}
Let $\bm Y$ be any vector target satisfying Assumption~\ref{as:context}, and put
$\bm\mu(\theta)=\E[\bm Y\mid\Theta=\theta]$ and
$\Sigma_Y(\theta)=\Cov(\bm Y\mid\Theta=\theta)$. Under squared-error loss,
the Bayes predictor is $\widehat{\bm Y}(C)=\E[\bm Y\mid C]$ and its risk is
\begin{equation}
\E\|\bm Y-\widehat{\bm Y}(C)\|_2^2
=\E\,\tr\!\left[\Cov\!\left(\bm\mu(\Theta)\mid C\right)\right]
+\E\,\tr\!\left[\Sigma_Y(\Theta)\right].
\label{eq:risk}
\end{equation}
\end{proposition}
\begin{proof}
Conditional independence gives
$\E[\bm Y\mid C]=\E[\bm\mu(\Theta)\mid C]$.
The law of total covariance, conditional on $C$, yields
$\Cov(\bm Y\mid C)=\Cov(\bm\mu(\Theta)\mid C)
+\E[\Sigma_Y(\Theta)\mid C]$.
Take traces and expectations.
\end{proof}

For the five-band log-power target in Eq.~\eqref{eq:bandpower}, if the
conditional distribution in Assumption~\ref{as:spectral} applies to the masked
patch, Eq.~\eqref{eq:logmean} gives
\begin{equation}
\mu_b(\theta)=\log s_b(\theta)+\kappa_b,
\qquad\kappa_b=\psi(m_b)-\log m_b,
\label{eq:mu}
\end{equation}
and Eq.~\eqref{eq:logvar} gives conditional variance $\psi_1(m_b)$.
Consequently, the state-independent part of the five-dimensional squared-error
risk is $\sum_b\psi_1(m_b)$ before any optional averaging over bands.
The other term is posterior uncertainty about the five log spectral levels that
remains even for the Bayes predictor with this fixed context. It is not reducible
merely by improving a predictor while keeping the same context. For any predictor
$f$, the actual excess risk is
\begin{equation}
R(f)-R^*=\E\|f(C)-\E[\bm Y\mid C]\|_2^2,
\end{equation}
where $R^*$ is the full Bayes risk in Eq.~\eqref{eq:risk}.
If one instead uses $\log(a+P_b)$, the conditional sampling term generally
depends on the state, as Eq.~\eqref{eq:logoffset} illustrates.

For waveform reconstruction, suppose that, conditional on a known $\theta$,
the masked sample vector $\bm x_M$ and an observed waveform context $\bm x_C$
are jointly zero-mean Gaussian with covariance blocks
$\Sigma_{MM},\Sigma_{MC},\Sigma_{CM},\Sigma_{CC}$, with $\Sigma_{CC}$ invertible.
Then
\begin{align}
\E[\bm x_M\mid\bm x_C,\theta]
&=\Sigma_{MC}\Sigma_{CC}^{-1}\bm x_C,\label{eq:wiener}\\
\Cov(\bm x_M\mid\bm x_C,\theta)
&=\Sigma_{MM}-\Sigma_{MC}\Sigma_{CC}^{-1}\Sigma_{CM}.
\label{eq:innovation}
\end{align}
The trace of Eq.~\eqref{eq:innovation}, divided by the number of masked samples,
is the conditional mean-squared \emph{waveform innovation} even when the state
is known. It includes unpredictable amplitude and phase, not phase alone.
Uncertainty about $\Theta$ can add further error to prediction from $C$.

Equations~\eqref{eq:risk}--\eqref{eq:innovation} motivate the hypothesis that
a band-power target can direct learning toward persistent spectral-state
information rather than the full waveform realization. They do \emph{not}
establish a universal ordering between the two objectives: their target
dimensions, units, loss normalizations, and Bayes-risk floors differ. Nor do
they prove better downstream classification, faster optimization, or greater
data efficiency. Such claims require controlled empirical comparisons.

\section{Where the model ceases to apply}

Several EEG properties lie outside this five-band account. Sleep spindles,
K-complexes, epileptiform events, and other transients may be distinguished by
shape and timing rather than average energy. Cross-channel coherence,
interregional phase relationships, phase--amplitude coupling, and fine spectral
structure are absent from five single-channel band powers. Eye movement,
muscle activity, and recording artifacts can contribute energy inside the same
bands as cortical rhythms. A state transition inside a patch violates the
locally persistent-state picture. Finally, finite windows limit frequency
resolution, especially for the lowest frequencies, as Table~\ref{tab:bins}
shows.

These are not minor technical exceptions. They delimit when spectral energy
is a plausible target for learning task-relevant structure and when it may
discard information that the downstream label actually needs.

\section{Conclusion}

Under an explicitly defined five-band Gaussian spectral model, the phases of
the observed Fourier coefficients are ancillary, the five classical band
powers are sufficient for the five spectral levels, and the log-power target
has state-independent conditional variance. A separate shared-state
assumption connects these properties to masked prediction through an exact
Bayes-risk decomposition. The analysis provides a principled reason to test
band-power prediction for labels tied to sustained EEG rhythms, while leaving
its comparative representation quality to experiment.

\end{document}
